% 农业上市公司画像：农业综合Ⅱ（共 2 家，按流通市值降序）
% 由 scripts/make_company_profiles.py 自动生成，请勿手工修改

\subsubsection{辉隆股份（002556）}
\label{co:002556}

\pfl{公司基本情况} 辉隆股份成立于 2011 年，安徽省供销社系统背景的农资流通龙头，形成化肥、农药等农资流通与精细化工“双主业”，并以 14 座现代农业综合服务中心为载体构建“耕种管收售”全周期农业服务体系。 公司于 2011-03-02 在深交所主板首发上市，注册资本 9.35 亿元，总股本 9.47 亿股。 截至 2026-09-24收盘价 5.08 元，流通市值 47.4 亿元，近一年-7.3\%，流通市值在三级行业“—”的 2 家公司中排第\zhnumber{1}位，近一年涨跌幅高于全样本中位数（-10.1\%）2.8 个百分点。 按 2026 年半年报披露的行业构成，收入占比前三位为农资产品 73.2\%；农副产品及其他 12.3\%；化工产品 8.1\%。 与 2021 年年报相比，第一大行业口径由“化肥”（62.1\%）变为“农资产品”（73.2\%）；两期披露的分类口径有调整，占比差异中同时包含分类因素。 发展过程中的关键节点包括：2020-02-25 因发行股份购买资产新增的 118,353,739 股（发行价格 5.13 元/股）上市流通前完成登记……（2020）；2023-07-28 发行 2023 年度第一期中期票据（乡村振兴）“23 辉隆农资 MTN001(乡村振兴)”……（2023）；2024-03-22 发行 2024 年度第一期中期票据（乡村振兴）“24 辉隆农资 MTN001(乡村振兴)”……（2024）。

\pfl{历史股价} 自 2021 年 1 月至 2026 年 9 月，股价由 8.11 元变为 5.08 元，累计 -37.4\%，区间最高 12.87 元（2021 年 12 月）、最低 4.17 元（2024 年 6 月）；当前价相当于区间最高价的 39.5\%，为区间最低价的 1.22 倍。 月度收益的年化波动率 30.5\%，低于全样本中位数 37.4\%，在三级行业“—”的 2 家公司中波动率排第\zhnumber{1}位。 从事件看，2024 年 6 月 至 2026 年 2 月的上涨（+61.2\%）与公司媒体报道控股股东解除 410 万股质押（与 2025 年末披露的质押 24,370,000 股口径存在差异……（2026-01-26）在时间上重合——股价的拐点多数能在公司自身的资本运作与风险事件上找到对应。

\begin{center}
\begin{minipage}{\textwidth}\centering
\begin{tikzpicture}
\begin{axis}[
  width=12.6cm, height=3.9cm, scale only axis,
  xmin=2020.922, xmax=2026.888, ymin=3.88, ymax=13.77,
  xtick={2021,2022,2023,2024,2025,2026},
  yearticks,
  yticklabel style={font=\scriptsize},
  ylabel={元/股}, ylabel style={font=\scriptsize},
  grid=major, grid style={LightGray, line width=0.3pt},
  axis line style={InkGray, line width=0.5pt},
  every axis plot/.append style={line width=0.9pt},
]
\addplot[AgriGreen] table[x=x,y=y]{data/clean/profiles/px_002556.dat};
\addplot[SoftRed, dashed, line width=0.7pt] table[x=x,y=y]{data/clean/profiles/pxzz_002556.dat};
\addplot[only marks, mark=*, mark size=1.1pt, SoftRed] table[x=x,y=y]{data/clean/profiles/pxzz_002556.dat};
\end{axis}
\end{tikzpicture}

\captionof{figure}[辉隆股份（002556）股价与周期骨架]{辉隆股份（002556）月度收盘价（元/股）与 ZigZag 周期骨架（阈值 25\%，圆点为周期转折点）}
\end{minipage}
\end{center}

\pfl{过去周期分析} 以 25\% 的 ZigZag 阈值划分，样本区间内股价形成 2 个主升段与 2 个主跌段。 最大主升段出现在 2024 年 6 月 至 2026 年 2 月，20 个月+61.2\%；最大主跌段为 2021 年 12 月 至 2024 年 6 月，30 个月-67.6\%。 最近一次转折出现在 2026 年 6 月（下跌段自 2026 年 2 月起，4 个月-33.3\%），当前处于该段的延续之中。 公司股价较申万农林牧渔指数提前约 23 个月见底，是板块中较早定价周期反转的公司之一。 农业综合类公司业务分散，估值更多取决于资产与转型预期而非单一品种价格，公司 2026 年 6 月末资产负债率 67.0\%，高于全样本中位数 52.2\%，其股价因此更多反映资金面与信用风险，而不是价格弹性本身。

\pfl{盈利情况} 2026 年上半年营业收入 85.47 亿元（同比 +3.3\%），归母净利润 1.13 亿元，同比 +1.5\%。 以年报为趋势对照，2023---2025 年营业收入由 178.33 亿元变为 151.32 亿元（两年年均复合增速 -7.9\%），归母净利润由 0.77 亿元变为 1.94 亿元。 按半年报口径，2026 年上半年的 ROE 为 3.1\%（未年化）、销售净利率 1.4\%、毛利率 6.2\%，ROE 在“—”2 家公司中按数值由高到低排第\zhnumber{1}位。 同期全样本半年 ROE 中位数 0.75\%，公司与之相差 2.33 个百分点（略好于中位数公司）。 从公开披露看，2026 年半年度计提资产减值准备 10,022.40 万元（项目主要为应收账款、其他应收款、存货、固定资产等）……；2026 年上半年经营活动产生的现金流量净额 -417,263,558.44 元……。

\begin{center}\small\setlength{\tabcolsep}{3.4pt}
\captionof{table}[辉隆股份（002556）关键财务指标]{辉隆股份（002556）关键财务指标（合并报表口径；两个半年报期均未年化；现金短债比 $=$ 货币资金 $\div$ 短期债务）}
\begin{adjustbox}{max width=\textwidth}
\begin{tabular}{@{}lrrrrr@{}}
\toprule
指标 & 2023 年 & 2024 年 & 2025 年 & 2025 年半年报 & 2026 年半年报 \\
\midrule
营业收入（亿元） & 178.33 & 156.52 & 151.32 & 82.77 & 85.47 \\
归母净利润（亿元） & 0.77 & 1.69 & 1.94 & 1.11 & 1.13 \\
毛利率（\%） & 4.8 & 5.8 & 6.3 & 5.8 & 6.2 \\
ROE（\%） & 2.0 & 4.6 & 5.2 & 3.0 & 3.1 \\
资产负债率（\%） & 66.2 & 64.6 & 64.9 & 67.9 & 67.0 \\
经营活动现金流净额（亿元） & -3.73 & -1.54 & 10.71 & -1.09 & -4.17 \\
货币资金（亿元） & 11.16 & 9.62 & 11.59 & 11.67 & 11.20 \\
有息负债合计（亿元） & 31.71 & 29.92 & 25.60 & 34.29 & 31.42 \\
现金短债比（倍） & 0.70 & 0.79 & 0.82 & 0.65 & 0.40 \\
\bottomrule
\end{tabular}
\end{adjustbox}
\end{center}

\begin{center}
\begin{minipage}{\textwidth}\centering
\begin{tikzpicture}
\begin{groupplot}[
  group style={group size=2 by 1, horizontal sep=0.60cm},
  width=6.85cm, height=4.60cm,
  tick label style={font=\tiny},
  title style={font=\scriptsize\heiti},
  yticklabel style={font=\tiny},
  ylabel style={font=\tiny},
  grid=major, grid style={LightGray, line width=0.3pt},
  axis line style={InkGray, line width=0.5pt},
  enlarge x limits={abs=0.8},
  legend style={font=\scriptsize, at={(0.5,-0.20)}, anchor=north,
                legend columns=2, draw=none, fill=none, column sep=6pt},
]
\nextgroupplot[
  ybar, bar width=4pt,
  ymin=-19.08, ymax=213.71,
  xtick={0,1,2,3,4,5.7},
  xticklabels={2021,2022,2023,2024,2025,2026H1},
  ylabel={亿元},
]
\addplot[fill=AgriGreen!75, draw=AgriGreen, bar shift=-2.6pt] table[x=i, y=rev]{data/clean/profiles/fin_002556.dat};
\addplot[fill=SoftRed!60, draw=SoftRed, bar shift=2.6pt] table[x=i, y=ni]{data/clean/profiles/fin_002556.dat};
\addplot[fill=AgriGreen!30, draw=AgriGreen, pattern=north east lines, bar shift=-2.6pt] table[x=x, y=rev]{data/clean/profiles/finh_002556.dat};
\addplot[fill=SoftRed!20, draw=SoftRed, pattern=north east lines, bar shift=2.6pt] table[x=x, y=ni]{data/clean/profiles/finh_002556.dat};
\legend{营业收入（年/半年）, 归母净利润（年/半年）}
\nextgroupplot[
  ymin=-5.4, ymax=73.7,
  xtick={0,1,2,3,4,5.7},
  xticklabels={2021,2022,2023,2024,2025,2026H1},
  ylabel={\%},
]
\addplot[AgriGold, mark=*, mark size=1.3pt] table[x=i, y=roe]{data/clean/profiles/fin_002556.dat};
\addplot[OilBlue, mark=square*, mark size=1.3pt] table[x=i, y=debt]{data/clean/profiles/fin_002556.dat};
\addplot[only marks, AgriGold, mark=*, mark size=2.0pt] table[x=x, y=roe]{data/clean/profiles/finh_002556.dat};
\addplot[only marks, OilBlue, mark=square*, mark size=2.0pt] table[x=x, y=debt]{data/clean/profiles/finh_002556.dat};
\legend{ROE, 资产负债率}
\end{groupplot}
\end{tikzpicture}
\begin{tikzpicture}
\begin{axis}[
  name=finbot,
  width=13.75cm, height=3.10cm, scale only axis,
  ymin=-3.17, ymax=35.52,
  xtick={0,1,2,3,4,5.7},
  xticklabels={2021,2022,2023,2024,2025,2026H1},
  tick label style={font=\tiny},
  yticklabel style={font=\tiny},
  ylabel={亿元}, ylabel style={font=\tiny},
  ybar, bar width=4pt,
  grid=major, grid style={LightGray, line width=0.3pt},
  axis line style={InkGray, line width=0.5pt},
  enlarge x limits={abs=0.8},
  legend style={font=\scriptsize, at={(0.5,-0.26)}, anchor=north,
                legend columns=2, draw=none, fill=none, column sep=10pt},
]
\addplot[fill=AgriGreen!55, draw=AgriGreen, bar shift=-3.0pt] table[x=i, y=cash]{data/clean/profiles/fin_002556.dat};
\addplot[fill=SoftRed!45, draw=SoftRed, bar shift=3.0pt] table[x=i, y=ibd]{data/clean/profiles/fin_002556.dat};
\addplot[fill=AgriGreen!25, draw=AgriGreen, pattern=north east lines, bar shift=-3.0pt] table[x=x, y=cash]{data/clean/profiles/finh_002556.dat};
\addplot[fill=SoftRed!20, draw=SoftRed, pattern=north east lines, bar shift=3.0pt] table[x=x, y=ibd]{data/clean/profiles/finh_002556.dat};
\legend{货币资金（左轴）, 有息负债（左轴）}
\end{axis}
\begin{axis}[
  at={(finbot.south east)}, anchor=south east,
  width=13.75cm, height=3.10cm, scale only axis,
  axis y line*=right, axis x line*=none, xtick=\empty,
  ymin=-8.49, ymax=17.41,
  tick label style={font=\tiny},
  yticklabel style={font=\tiny}, scaled y ticks=false,
  legend style={font=\scriptsize, at={(0.5,-0.44)}, anchor=north,
                legend columns=1, draw=none, fill=none},
]
\addplot[OilBlue, mark=*, mark size=2.2pt, line width=1.3pt] table[x=i, y=ocf]{data/clean/profiles/fin_002556.dat};
\addplot[only marks, OilBlue, mark=*, mark size=3.2pt] table[x=x, y=ocf]{data/clean/profiles/finh_002556.dat};
\legend{经营活动现金流净额（右轴，亿元，蓝点）}
\end{axis}
\end{tikzpicture}

\captionof{figure}[辉隆股份（002556）近年主要财务指标]{辉隆股份（002556）近年主要财务指标（左上：营业收入与归母净利润；右上：ROE 与资产负债率；下：货币资金、有息负债为左轴柱状，经营活动现金流净额为其右轴的蓝点折线。
2021---2025 年为年报口径，2026H1 为 2026 年半年报/半年末，与其前一年报之间留出间隔、以斜纹或空心点区分；半年报数据未年化）}
\end{minipage}
\end{center}

\pfl{财务分析} 2026 年 6 月末资产负债率 67.0\%，较 2025 年末 上升 2.1 个百分点（样本中位数 52.2\%，所处三级行业内按由低到高排第\zhnumber{1}位）。 2026 年 6 月末货币资金 11.20 亿元、短期债务（短期借款与一年内到期非流动负债）27.68 亿元、长期债务 3.74 亿元，有息负债合计 31.42 亿元，现金短债比 0.40 倍——覆盖偏紧。 净有息负债 20.22 亿元，相当于其 2026 年上半年营业收入的 24\%。 流动比率 1.01、速动比率 0.55，短期指标尚可。 2026 年上半年经营活动现金流净额 -4.17 亿元，净利润为正但经营现金流为负，利润的现金含量偏低。 公司披露的财务与合规风险包括：受限货币资金 384,448,852.39 元，受限类型为保证，用于开具银行承兑汇票、信用证、保函、借款……；媒体报道控股股东解除 410 万股质押（与 2025 年末披露的质押 24,370,000 股口径存在差异……。

\pfl{重大战略转型} 近三年公司的战略动作集中在以下几个方面：2026-02-04 媒体报道公司全资孙公司中成科技位于硫铁矿带；参股企业（中盟磷业）白岩磷矿储量可达 9,500 万吨；2026-04-17 披露《关于“质量回报双提升”行动方案的公告》（公告编号 2026-021），2026-08-13 披露进展：注销回购股份 11,577,228 股……；2026-06-30 聚焦“农资+精细化工”双主业提质增效，从资源端、产品端和市场端夯实主业基石，同时推进科研应用新赛道转型；2026-06-30 农服板块围绕种肥药一体化服务延伸链条，对接优质种质资源、落地绿色农资示范项目、新增秧苗培育业务。 公告口径上，近三年（2023 年 9 月---2026 年 9 月）公司披露重大重组与并购类公告 0 条、对外投资与转型类公告 1 条、回购与股权激励类公告 15 条，合计公告 341 条。 主营结构的口径变化已经落到报表上：2021 年年报的第一大行业为“化肥”（62.1\%），2026 年半年报为“农资产品”（73.2\%）；公司在这两期的分部划分方式有过调整。 综合判断，公司在报告期内有明确的外延扩张或业务结构调整动作，转型方向与融资安排之间存在直接关联。

\begin{center}\small\setlength{\tabcolsep}{4pt}
\captionof{table}[辉隆股份（002556）历史融资与资本运作]{辉隆股份（002556）历史融资与资本运作（据公司公告与公开报道整理；金额为披露值，含首次公开发行、再融资、债券、银行借款与重整投资等）}
\begin{adjustbox}{max width=\textwidth}
\begin{tabular}{@{}p{1.45cm}p{2.55cm}p{4.05cm}p{4.95cm}@{}}
\toprule
时间 & 方式 & 规模 & 用途与说明 \\
\midrule
2011-02 & IPO（深圳证券交易所中小板） & 发行 3,\allowbreak{}750 万股，\allowbreak{}发行价 38 元/\allowbreak{}股，\allowbreak{}募集资金总额 14.06 亿元；招股说明书披露…… & 配送中心建设项目和信息化系统建设项目；经证监许可[2011]197 号文核准…… \\
2018-2019 /\allowbreak{} 2020-02 & 发行股份、可转换公司债券及支付现金购买资产（关联交易） & 标的为海华科技 100\% 股权，\allowbreak{}总对价 8.28 亿元：现金 1.81 亿元、\allowbreak{}定向可转债 4…… & 收购海华科技 100\% 股权（交易对方为辉隆投资、蚌埠隆海及解凤贤、解凤苗…… \\
2020-04 & 非公开发行股份及可转换公司债券募集配套资金 & 募集配套资金 6.44 亿元，\allowbreak{}其中非公开发行股份募集 1.30 亿元、\allowbreak{}非公开发行可转换公司债…… & 向交易对方支付现金对价、海华科技项目建设及支付本次交易相关费用…… \\
2023-07-28 & 中期票据（乡村振兴） & 4.00 亿元，\allowbreak{}票面利率 3.50\%，\allowbreak{}到期日 2028-08-01（债券代码 1023818…… & 2025 年度报告披露期末余额 405,833,333.33 元…… \\
2024-03-22 & 中期票据（乡村振兴） & 4.00 亿元，\allowbreak{}票面利率 2.80\%，\allowbreak{}到期日 2027-03-26（债券代码 1024811…… & 2025 年度报告披露期末余额 408,788,888.89 元…… \\
2026-04-17 & 终止部分超募资金投资项目 & 剩余超募资金 4,\allowbreak{}343.59 万元 & 终止首发超募资金投资项目“40 万吨/年新型肥料项目”…… \\
\bottomrule
\end{tabular}
\end{adjustbox}
\end{center}

\pfl{历史融投资情况} 从现金流量表看，2025 年公司取得借款收到的现金 18.70 亿元、偿还债务支付的现金 25.51 亿元、吸收权益性投资 0.01 亿元，筹资活动现金流净额 -10.40 亿元（2023 年为取得借款 36.24 亿元、偿还债务 28.20 亿元）。 2026 年上半年筹资活动现金流净额为 0.00 亿元（取得借款 14.36 亿元、偿还债务 13.04 亿元，未年化），仍处于净融入状态。 2025 年分配股利、利润或偿付利息支付的现金合计 3.18 亿元。 公司披露的股本与债务融资脉络包括：IPO（深圳证券交易所中小板）、中期票据（乡村振兴）、发行股份、可转换公司债券及支付现金购买资产（关联交易）、终止部分超募资金投资项目、非公开发行股份及可转换公司债券募集配套资金。 其中规模较大的两笔为：2026-04-17终止部分超募资金投资项目（剩余超募资金 4,343.59 万元）；2020-04非公开发行股份及可转换公司债券募集配套资金（募集配套资金 64,406.52 万元，其中非公开发行股份募集 13,000 万元……）。 股本方面，近三年披露股本变动 9 次（其他 3 次、限售股份上市 2 次、股份回购 1 次）；总股本由 2021 年末的 9.54 亿股变为 9.47 亿股（-0.78\%）。 分红方面，近三年现金分红 5 次、累计每股派现 0.85 元（最近一次 2025-12-31）——在利润承压的阶段仍维持了分红。

\pfl{未来潜在融资需求分析} 2026 年 6 月末货币资金 11.20 亿元，而短期债务 27.68 亿元（现金短债比 0.40 倍），2026 年上半年经营现金流净额 -4.17 亿元。按“短期债务减去账面货币资金”的滚动偿付口径，静态缺口约 16.47 亿元；若再计入上半年亏损的年化额（0.00 亿元），维持一年正常经营所需的外部资金约 16.47---16.47 亿元。 公司 2026 年 6 月末资产负债率 67.0\%、2026 年上半年 ROE 3.1\%（未年化），资产负债表尚有余量，融资需求不是“补血型”的刚性需求，而是围绕并购整合、项目投入与营运资金的主动性需求。 公开披露的下一步融资线索包括：2026-05 2026-05-07 披露变更回购股份用途并注销、减少注册资本及修订《公司章程》……；2026-06-30 披露对控股公司提供担保及募集资金存放与实际使用情况专项报告；货币资金 1,030,458,766.52 元；2026-08-07 披露《关于 2023 年度第一期中期票据(乡村振兴)付息及回售兑付完成的公告》，4 亿元中期票据进入回售兑付环节。




\subsubsection{大禹节水（300021）}
\label{co:300021}

\pfl{公司基本情况} 大禹节水成立于 2009 年，国内节水灌溉与农田水利领域主要上市公司，业务覆盖节水材料装备制造、灌区 EPC/EPCO 总承包与农水设计服务，正向“设计引领+全产业链”模式转型。 公司于 2009-10-30 在深交所创业板首发上市，注册资本 10.22 亿元，总股本 8.54 亿股。 截至 2026-09-24收盘价 3.82 元，流通市值 33.4 亿元，近一年-22.7\%，流通市值在三级行业“—”的 2 家公司中排第\zhnumber{2}位，近一年涨跌幅低于全样本中位数（-10.1\%）12.6 个百分点。 按 2025 年年报披露的行业构成，收入占比前三位为智慧水利行业 100.0\%。 与 2021 年年报相比，第一大行业口径由“节水灌溉行业”（81.0\%）变为“智慧水利行业”（100.0\%）；两期披露的分类口径有调整，占比差异中同时包含分类因素。 发展过程中的关键节点包括：2010-04-14 召开 2009 年年度股东大会，审议通过以 2009 年末总股本 69,650,000 股为基数……（2010）；2020-07-28 公开发行 638 万张可转换公司债券（每张面值 100 元，发行总额 6.38 亿元）……（2020）；2022-02-24 以简易程序向特定对象发行股票 58,593,750 股（发行价 5.12 元/股）募集资金到账……（2022）。

\pfl{历史股价} 以 2021 年 1 月为起点、2026 年 9 月为终点，股价由 4.76 元变为 3.82 元，累计 -19.7\%，区间最高 6.32 元（2025 年 7 月）、最低 3.25 元（2024 年 8 月）；当前价相当于区间最高价的 60.4\%，为区间最低价的 1.18 倍。 月度收益的年化波动率 27.2\%，低于全样本中位数 37.4\%，在三级行业“—”的 2 家公司中波动率排第\zhnumber{2}位。 从事件看，2024 年 8 月 至 2025 年 7 月的上涨（+94.5\%）与公司可转债募集资金：本报告期投入 8,418.05 万元……（2025-06-30）在时间上重合；2021 年 12 月 至 2024 年 8 月的下跌（-44.9\%）与公司股东大会通过变更部分募集资金用途，将可转债募投项目尚未使用募集资金及专户存款……（2024-02-05）在时间上重合——股价的拐点多数能在公司自身的资本运作与风险事件上找到对应。

\begin{center}
\begin{minipage}{\textwidth}\centering
\begin{tikzpicture}
\begin{axis}[
  width=12.6cm, height=3.9cm, scale only axis,
  xmin=2020.922, xmax=2026.888, ymin=3.02, ymax=6.76,
  xtick={2021,2022,2023,2024,2025,2026},
  yearticks,
  yticklabel style={font=\scriptsize},
  ylabel={元/股}, ylabel style={font=\scriptsize},
  grid=major, grid style={LightGray, line width=0.3pt},
  axis line style={InkGray, line width=0.5pt},
  every axis plot/.append style={line width=0.9pt},
]
\addplot[AgriGreen] table[x=x,y=y]{data/clean/profiles/px_300021.dat};
\addplot[SoftRed, dashed, line width=0.7pt] table[x=x,y=y]{data/clean/profiles/pxzz_300021.dat};
\addplot[only marks, mark=*, mark size=1.1pt, SoftRed] table[x=x,y=y]{data/clean/profiles/pxzz_300021.dat};
\end{axis}
\end{tikzpicture}

\captionof{figure}[大禹节水（300021）股价与周期骨架]{大禹节水（300021）月度收盘价（元/股）与 ZigZag 周期骨架（阈值 25\%，圆点为周期转折点）}
\end{minipage}
\end{center}

\pfl{过去周期分析} 以 25\% 的 ZigZag 阈值划分，样本区间内股价形成 2 个主升段与 2 个主跌段。 最大主跌段为 2021 年 12 月 至 2024 年 8 月，32 个月-44.9\%；最大主升段出现在 2024 年 8 月 至 2025 年 7 月，11 个月+94.5\%。 最近一次转折出现在 2026 年 9 月（下跌段自 2025 年 7 月起，14 个月-39.6\%），当前处于该段的延续之中。 公司股价较申万农林牧渔指数提前约 21 个月见底，是板块中较早定价周期反转的公司之一。 农业综合类公司业务分散，估值更多取决于资产与转型预期而非单一品种价格，公司 2026 年 6 月末资产负债率 71.7\%，高于全样本中位数 52.2\%，其股价因此更多反映资金面与信用风险，而不是价格弹性本身。

\pfl{盈利情况} 2026 年上半年营业收入 15.86 亿元（同比 +24.3\%），归母净利润 -0.01 亿元，由上年同期的盈利转为亏损。 以年报为趋势对照，2023---2025 年营业收入由 34.53 亿元变为 37.61 亿元（两年年均复合增速 +4.4\%），归母净利润由 0.50 亿元变为 0.49 亿元。 按半年报口径，2026 年上半年的 ROE 为 -0.0\%（未年化）、销售净利率 -0.4\%、毛利率 20.6\%，ROE 在“—”2 家公司中按数值由高到低排第\zhnumber{2}位。 同期全样本半年 ROE 中位数 0.75\%，公司与之相差 0.79 个百分点（弱于中位数公司）。 从公开披露看，转让所持慧图科技 28.5\% 股权给嘉兴乾瞻衡远创业投资合伙企业（有限合伙），交易价格 19,237.50 万元……；2025 年度营业收入 376,124.83 万元，同比下降 14.10\%；归属于上市公司股东的净利润 4,890.01 万元……。

\begin{center}\small\setlength{\tabcolsep}{3.4pt}
\captionof{table}[大禹节水（300021）关键财务指标]{大禹节水（300021）关键财务指标（合并报表口径；两个半年报期均未年化；现金短债比 $=$ 货币资金 $\div$ 短期债务）}
\begin{adjustbox}{max width=\textwidth}
\begin{tabular}{@{}lrrrrr@{}}
\toprule
指标 & 2023 年 & 2024 年 & 2025 年 & 2025 年半年报 & 2026 年半年报 \\
\midrule
营业收入（亿元） & 34.53 & 43.79 & 37.61 & 12.76 & 15.86 \\
归母净利润（亿元） & 0.50 & 0.81 & 0.49 & 0.13 & -0.01 \\
毛利率（\%） & 24.3 & 17.8 & 21.4 & 23.3 & 20.6 \\
ROE（\%） & 2.5 & 3.9 & 2.2 & 0.6 & -0.0 \\
资产负债率（\%） & 71.2 & 75.1 & 70.6 & 73.6 & 71.7 \\
经营活动现金流净额（亿元） & -0.94 & 5.97 & -2.76 & -5.49 & -5.69 \\
货币资金（亿元） & 11.51 & 16.88 & 14.69 & 12.02 & 12.85 \\
有息负债合计（亿元） & 25.14 & 27.10 & 25.96 & 29.00 & 29.12 \\
现金短债比（倍） & 1.14 & 1.22 & 0.98 & 0.72 & 0.70 \\
\bottomrule
\end{tabular}
\end{adjustbox}
\end{center}

\begin{center}
\begin{minipage}{\textwidth}\centering
\begin{tikzpicture}
\begin{groupplot}[
  group style={group size=2 by 1, horizontal sep=0.60cm},
  width=6.85cm, height=4.60cm,
  tick label style={font=\tiny},
  title style={font=\scriptsize\heiti},
  yticklabel style={font=\tiny},
  ylabel style={font=\tiny},
  grid=major, grid style={LightGray, line width=0.3pt},
  axis line style={InkGray, line width=0.5pt},
  enlarge x limits={abs=0.8},
  legend style={font=\scriptsize, at={(0.5,-0.20)}, anchor=north,
                legend columns=2, draw=none, fill=none, column sep=6pt},
]
\nextgroupplot[
  ybar, bar width=4pt,
  ymin=-4.39, ymax=49.04,
  xtick={0,1,2,3,4,5.7},
  xticklabels={2021,2022,2023,2024,2025,2026H1},
  ylabel={亿元},
]
\addplot[fill=AgriGreen!75, draw=AgriGreen, bar shift=-2.6pt] table[x=i, y=rev]{data/clean/profiles/fin_300021.dat};
\addplot[fill=SoftRed!60, draw=SoftRed, bar shift=2.6pt] table[x=i, y=ni]{data/clean/profiles/fin_300021.dat};
\addplot[fill=AgriGreen!30, draw=AgriGreen, pattern=north east lines, bar shift=-2.6pt] table[x=x, y=rev]{data/clean/profiles/finh_300021.dat};
\addplot[fill=SoftRed!20, draw=SoftRed, pattern=north east lines, bar shift=2.6pt] table[x=x, y=ni]{data/clean/profiles/finh_300021.dat};
\legend{营业收入（年/半年）, 归母净利润（年/半年）}
\nextgroupplot[
  ymin=-6.0, ymax=82.6,
  xtick={0,1,2,3,4,5.7},
  xticklabels={2021,2022,2023,2024,2025,2026H1},
  ylabel={\%},
]
\addplot[AgriGold, mark=*, mark size=1.3pt] table[x=i, y=roe]{data/clean/profiles/fin_300021.dat};
\addplot[OilBlue, mark=square*, mark size=1.3pt] table[x=i, y=debt]{data/clean/profiles/fin_300021.dat};
\addplot[only marks, AgriGold, mark=*, mark size=2.0pt] table[x=x, y=roe]{data/clean/profiles/finh_300021.dat};
\addplot[only marks, OilBlue, mark=square*, mark size=2.0pt] table[x=x, y=debt]{data/clean/profiles/finh_300021.dat};
\legend{ROE, 资产负债率}
\end{groupplot}
\end{tikzpicture}
\begin{tikzpicture}
\begin{axis}[
  name=finbot,
  width=13.75cm, height=3.10cm, scale only axis,
  ymin=-2.91, ymax=32.62,
  xtick={0,1,2,3,4,5.7},
  xticklabels={2021,2022,2023,2024,2025,2026H1},
  tick label style={font=\tiny},
  yticklabel style={font=\tiny},
  ylabel={亿元}, ylabel style={font=\tiny},
  ybar, bar width=4pt,
  grid=major, grid style={LightGray, line width=0.3pt},
  axis line style={InkGray, line width=0.5pt},
  enlarge x limits={abs=0.8},
  legend style={font=\scriptsize, at={(0.5,-0.26)}, anchor=north,
                legend columns=2, draw=none, fill=none, column sep=10pt},
]
\addplot[fill=AgriGreen!55, draw=AgriGreen, bar shift=-3.0pt] table[x=i, y=cash]{data/clean/profiles/fin_300021.dat};
\addplot[fill=SoftRed!45, draw=SoftRed, bar shift=3.0pt] table[x=i, y=ibd]{data/clean/profiles/fin_300021.dat};
\addplot[fill=AgriGreen!25, draw=AgriGreen, pattern=north east lines, bar shift=-3.0pt] table[x=x, y=cash]{data/clean/profiles/finh_300021.dat};
\addplot[fill=SoftRed!20, draw=SoftRed, pattern=north east lines, bar shift=3.0pt] table[x=x, y=ibd]{data/clean/profiles/finh_300021.dat};
\legend{货币资金（左轴）, 有息负债（左轴）}
\end{axis}
\begin{axis}[
  at={(finbot.south east)}, anchor=south east,
  width=13.75cm, height=3.10cm, scale only axis,
  axis y line*=right, axis x line*=none, xtick=\empty,
  ymin=-8.72, ymax=9.46,
  tick label style={font=\tiny},
  yticklabel style={font=\tiny}, scaled y ticks=false,
  legend style={font=\scriptsize, at={(0.5,-0.44)}, anchor=north,
                legend columns=1, draw=none, fill=none},
]
\addplot[OilBlue, mark=*, mark size=2.2pt, line width=1.3pt] table[x=i, y=ocf]{data/clean/profiles/fin_300021.dat};
\addplot[only marks, OilBlue, mark=*, mark size=3.2pt] table[x=x, y=ocf]{data/clean/profiles/finh_300021.dat};
\legend{经营活动现金流净额（右轴，亿元，蓝点）}
\end{axis}
\end{tikzpicture}

\captionof{figure}[大禹节水（300021）近年主要财务指标]{大禹节水（300021）近年主要财务指标（左上：营业收入与归母净利润；右上：ROE 与资产负债率；下：货币资金、有息负债为左轴柱状，经营活动现金流净额为其右轴的蓝点折线。
2021---2025 年为年报口径，2026H1 为 2026 年半年报/半年末，与其前一年报之间留出间隔、以斜纹或空心点区分；半年报数据未年化）}
\end{minipage}
\end{center}

\pfl{财务分析} 2026 年 6 月末资产负债率 71.7\%，较 2025 年末 上升 1.1 个百分点（样本中位数 52.2\%，所处三级行业内按由低到高排第\zhnumber{2}位）。 2026 年 6 月末货币资金 12.85 亿元、短期债务（短期借款与一年内到期非流动负债）18.48 亿元、长期债务 10.64 亿元，有息负债合计 29.12 亿元，现金短债比 0.70 倍——覆盖偏紧。 净有息负债 16.27 亿元，相当于其 2026 年上半年营业收入的 103\%。 流动比率 1.25、速动比率 1.20，短期指标尚可。 2026 年上半年经营活动现金流净额 -5.69 亿元，经营现金流与利润同时为负，日常运营对债务与股东投入的依赖度较高。 2026 年 6 月末在建工程 0.61 亿元，较上年末增加 0.54 亿元，仍有在建投入。 2026 年 6 月末商誉 1.34 亿元，占归母净资产的 4.6\%，是净资产中最脆弱的一块。

\pfl{重大战略转型} 公司在报告期内的战略动作可归为 4 项：2026-03 转让所持慧图科技 28.5\% 股权给嘉兴乾瞻衡远创业投资合伙企业（有限合伙），交易价格 19,237.50 万元……；2026-06-30 依托“设计引领+全产业链”协同模式，设计业务成为第一增长引擎（农水设计服务收入同比 +306.84\%）……；2026-06-30 在生态治理与再生水新赛道取得突破，取得新疆阿克苏地区塔里木河干流盖孜库木片区生态综合治理项目……；2026-06-30 全面落地全域一体化营销机制（“大营销”协同机制），打通各业务板块资源壁垒，统筹全国区域项目联合开发，持续储备大中型灌区、流域治理……。 公告口径上，近三年（2023 年 9 月---2026 年 9 月）公司披露重大重组与并购类公告 9 条、对外投资与转型类公告 35 条、回购与股权激励类公告 42 条，合计公告 633 条。 主营结构的口径变化已经落到报表上：2021 年年报的第一大行业为“节水灌溉行业”（81.0\%），2025 年年报为“智慧水利行业”（100.0\%）；公司在这两期的分部划分方式有过调整。 综合判断，公司在报告期内有明确的外延扩张或业务结构调整动作，转型方向与融资安排之间存在直接关联。

\begin{center}\small\setlength{\tabcolsep}{4pt}
\captionof{table}[大禹节水（300021）历史融资与资本运作]{大禹节水（300021）历史融资与资本运作（据公司公告与公开报道整理；金额为披露值，含首次公开发行、再融资、债券、银行借款与重整投资等）}
\begin{adjustbox}{max width=\textwidth}
\begin{tabular}{@{}p{1.45cm}p{2.55cm}p{4.05cm}p{4.95cm}@{}}
\toprule
时间 & 方式 & 规模 & 用途与说明 \\
\midrule
2009-10 & IPO（深圳证券交易所创业板） & 发行 1,\allowbreak{}800 万股，\allowbreak{}发行价 14 元/\allowbreak{}股，\allowbreak{}预计募集资金总额 2.52 亿元、\allowbreak{}净额 2.2…… & 在酒泉总部及全资子公司新疆公司、武威公司建设滴灌管（带）及配套节水器材生产线技改扩建项目…… \\
2020-07-28 & 公开发行可转换公司债券“大禹转债” & 638 万张、\allowbreak{}每张面值 100 元，\allowbreak{}发行总额 6.38 亿元，\allowbreak{}期限 6 年（2020-07-…… & 高端节水灌溉产品智能工厂建设项目（拟配备以色列先进压力补偿式滴灌生产线及核心智能制造设备…… \\
2022-02-24 & 以简易程序向特定对象发行股票 & 58,\allowbreak{}593,\allowbreak{}750 股，\allowbreak{}发行价 5 元/\allowbreak{}股，\allowbreak{}募集资金总额 3.00 亿元，\allowbreak{}扣除发行费用 9…… & 未查得（募投项目及使用情况见年度报告募集资金专项说明）…… \\
2025-08-21 & 可转债提前赎回 & 赎回价格 100 元/\allowbreak{}张（含税）；公司披露本年度实际赎回 10,\allowbreak{}724.00 张大禹转债 & 终止“大禹转债”存续并在深交所摘牌；本年度因可转债转股增加股本 17,234.74 万股…… \\
2026-07 & 2025 年度利润分配 & 以 1,\allowbreak{}022,\allowbreak{}367,\allowbreak{}469 股为基数，\allowbreak{}每 10 股派发现金红利 0 元（含税），\allowbreak{}合计 2…… & 股东回报；2026 年 7 月完成派发 \\
\bottomrule
\end{tabular}
\end{adjustbox}
\end{center}

\pfl{历史融投资情况} 从现金流量表看，2025 年公司取得借款收到的现金 21.76 亿元、偿还债务支付的现金 16.85 亿元、吸收权益性投资 0.08 亿元，筹资活动现金流净额 2.39 亿元（2023 年为取得借款 14.48 亿元、偿还债务 9.61 亿元）。 2026 年上半年筹资活动现金流净额为 2.82 亿元（取得借款 10.47 亿元、偿还债务 7.10 亿元，未年化），仍处于净融入状态。 2025 年分配股利、利润或偿付利息支付的现金合计 1.46 亿元。 公司披露的股本与债务融资脉络包括：2025 年度利润分配、IPO（深圳证券交易所创业板）、以简易程序向特定对象发行股票、公开发行可转换公司债券“大禹转债”、可转债提前赎回。 其中规模较大的两笔为：2020-07-28公开发行可转换公司债券“大禹转债”（638 万张、每张面值 100 元，发行总额 6.38 亿元……）；2009-10IPO（深圳证券交易所创业板）（发行 1,800 万股，发行价 14.00 元/股，预计募集资金总额 25,200 万元……）。 股本方面，近三年披露股本变动 17 次（可转债转股 11 次、股份回购 4 次、其他 1 次）；总股本由 2021 年末的 8.01 亿股变为 8.54 亿股（+6.59\%）。 分红方面，近三年现金分红 4 次、累计每股派现 0.21 元（最近一次 2025-12-31）——分红规模与利润的下滑幅度相比仍然有限。 近三年“再融资”类公告 18 条，涉及定向增发、可转债。

\pfl{未来潜在融资需求分析} 2026 年 6 月末货币资金 12.85 亿元，而短期债务 18.48 亿元（现金短债比 0.70 倍），2026 年上半年经营现金流净额 -5.69 亿元。按“短期债务减去账面货币资金”的滚动偿付口径，静态缺口约 5.62 亿元；若再计入上半年亏损的年化额（0.02 亿元），维持一年正常经营所需的外部资金约 5.62---5.65 亿元。 按第 \ref{sec:financing-need} 节的三档划分，公司属于第一档“补血型”：2026 年上半年归母净利润 -0.01 亿元、6 月末资产负债率 71.7\%。 可行的融资路径按现实性排序为：定向增发（需股价配合，且控股股东需放弃优先认购或同步增资）、出售非核心资产与参股股权、债务重组或展期（与银行和债券持有人协商）。 公开披露的下一步融资线索包括：2025-06-30 可转债募集资金：本报告期投入 8,418.05 万元，截至 2025-06-30 累计使用 33,062.88 万元……；2025-08-26 公司提出高标准农田建设由整改前单体项目拆分招标转为整县整装推进，将显著提升单个订单体量，对具有规模优势的上市公司参与此类项目更有利……；2025-12-31 简易程序发行募集资金：本报告期投入 4,375.47 万元，截至 2025-12-31 累计使用 20,635.73 万元……。 资金需求还要叠加在建工程：期末在建工程 0.61 亿元、2026 年上半年购建固定资产等支出 0.45 亿元（未年化），这部分支出在周期底部通常会被主动放缓。




